\section{Proof details}
\label{app:proofs}

\subsection{Triangular affine pullback}

For \(\phi(z)=a+qz\),
\[
 U_{\phi,0}z^k=(a+qz)^k
 =\sum_{j=0}^k\binom{k}{j}a^{k-j}q^jz^j.
\]
Thus the finite monomial matrix is upper triangular by degree and has
diagonal \(1,q,\ldots,q^N\).  On one-forms,
\[
 U_{\phi,1}(z^kdz)=q(a+qz)^kdz,
\]
whose diagonal is \(q,q^2,\ldots,q^N\) on
\(P_{N-1}dz\).  This proves directly that
\begin{align*}
 \Tr(wU_{\phi,0}|P_N)^r
 &=w^r\sum_{m=0}^Nq^{rm},\\
 \Tr(wU_{\phi,1}|P_{N-1}dz)^r
 &=w^r\sum_{m=1}^Nq^{rm},
\end{align*}
so their difference is \(w^r\) for every \(N\) and \(r\).

To check the chain identity on a monomial, compute
\begin{align*}
 dU_{\phi,0}z^k
 &=kq(a+qz)^{k-1}dz,\\
 U_{\phi,1}d z^k
 &=U_{\phi,1}(kz^{k-1}dz)
 =kq(a+qz)^{k-1}dz.
\end{align*}
This is \(DM_0=M_1D\) in matrix form.

\subsection{Finite cohomological determinant identity}

Choose the invariant constants \(C\subset P_N\) and the quotient
\(P_N/C\).  Differentiation induces an isomorphism
\[
        \overline d:P_N/C\longrightarrow P_{N-1}dz.
\]
The chain identity makes the induced zero-form quotient action similar to
the one-form action.  In a basis adapted to \(C\subset P_N\), the weighted
zero-form matrix is block upper triangular with the scalar \(w\) on \(C\)
and a block similar to \(wM_1\) on the quotient.  Therefore
\[
 \det(I-zwM_0)=(1-zw)\det(I-zwM_1).
\]

For a weighted sum \(L_k=\sum_jw_jM_{j,k}\), constants carry
\(S_w=\sum_jw_j\), and the same quotient similarity yields
\[
 \det(I-zL_0)=(1-zS_w)\det(I-zL_1).
\]
No commutativity among the branch matrices is used.

\subsection{Tensor moment obstruction in full detail}

For trace-class \(B\), the Fredholm determinant has the local expansion
\[
 \log\det(I-tB)=-\sum_{r\ge1}\frac{t^r}{r}\Tr B^r.
\]
If \(\Tr B^r=1-q^r\), then
\begin{align*}
 \log\det(I-tB)
 &=-\sum_{r\ge1}\frac{t^r}{r}
   +\sum_{r\ge1}\frac{(qt)^r}{r}\\
 &=\log(1-t)-\log(1-qt).
\end{align*}
Exponentiation gives \((1-t)/(1-qt)\) on a neighborhood of zero.  By
analytic continuation it would equal the entire Fredholm determinant away
from its alleged pole, which is impossible.  The argument applies equally
to a finite matrix.  It uses the full moment sequence; a search over small
matrices is only a regression test.

\subsection{Absolute fixed-word summation}

For the code inventory \(n\ge2\), one has \(|q_n|\le2^{-3}\).  Hence for
every nonempty word \(\alpha\),
\[
 \left|\frac1{1-q_\alpha}\right|
 \le\frac1{1-2^{-3}}.
\]
If \(\sum_j|w_j|<\infty\), then at power \(r\)
\[
 \sum_{\alpha\in J^r}|w_\alpha|
 =\left(\sum_j|w_j|\right)^r<\infty.
\]
The degree-zero and degree-one fixed-word traces are therefore absolutely
summable.  Termwise subtraction is justified and yields
\[
 \sum_{\alpha\in J^r}
 w_\alpha\left(\frac1{1-q_\alpha}
 -\frac{q_\alpha}{1-q_\alpha}\right)
 =\sum_{\alpha\in J^r}w_\alpha=S_w^r.
\]

\subsection{Primitive-necklace identity}

Let \(A_r=(\sum_jw_j)^r\).  The shared connected trace logarithm is
\[
 \sum_{r\ge1}\frac{z^r}{r}A_r
 =-\log\left(1-z\sum_jw_j\right).
\]
Every word has a unique representation as a positive power of a primitive
cyclic word together with a choice of starting point.  Grouping the word sum
by this primitive root gives
\[
 -\log\left(1-z\sum_jw_j\right)
 =-\sum_{[\alpha]\ \mathrm{primitive}}
   \log(1-z^{|\alpha|}w_\alpha),
\]
which exponentiates to \eqref{eq:necklace-product}.  Pure one-letter roots
generate the repeated terms \(w_j^r\); roots using multiple labels generate
the mixed ledger.

\subsection{Direct-sum determinant and trace class}

For a direct sum \(T=\bigoplus_nT_n\), the sufficient trace-class condition
\(\sum_n\lVert T_n\rVert_1<\infty\) gives
\[
        \det(I-zT)=\prod_n\det(I-zT_n)
\]
locally uniformly in \(z\).  Here the common compact-containment estimate
gives \(\lVert U_{\phi_n,k}\rVert_1\le C_k\), uniformly in \(n\).  Thus
\[
 \sum_{n\in S}\lVert n^{-s}U_{\phi_n,k}\rVert_1
 \le C_k\sum_{n\ge2}n^{-\Re s}<\infty
\]
for \(\Re s>1\).  Componentwise application of
\eqref{eq:local-det-identity} proves \eqref{eq:disjoint-det}.

\subsection{Prime cohomology boundary}

The diagonal cohomology operator \(C_{\Pp,s}\) is normal with singular
values \(p^{-\Re s}\).  Hence
\[
        \lVert C_{\Pp,s}\rVert_1=\sum_pp^{-\Re s}.
\]
Euler's divergence of \(\sum_pp^{-1}\) gives the boundary at \(1\).
This is an operator statement, independent of any meromorphic continuation
of the scalar determinant ratio.
